% Ch10 WS3 -- vapour pressure, heating curves, phase-change energy + FRQ.
\documentclass{apchem-worksheet}

\apcunitword{Chapter}
\apcunit{10}
\apcunittitle{Liquids and Solids}
\apcdoctitle{Worksheet 3 \textbullet\ Vapour Pressure and Phase Changes}
\apcsubtitle{Zumdahl \S10.8 \textbullet\ also CED 6.5 --- Unit 6 borrows this topic}

\begin{document}
\wsheader[Water data for all problems: $c_{\text{ice}} = 2.09$,
$c_{\text{water}} = 4.184$, $c_{\text{steam}} = 2.0$
\si{\joule\per\gram\per\celsius}; $\Delta H_{\text{fus}} =
\SI{6.02}{\kilo\joule\per\mole}$;
$\Delta H_{\text{vap}} = \SI{40.7}{\kilo\joule\per\mole}$;
$M = \SI{18.02}{\gram\per\mole}$.]

\problem[]{Vapour pressure concepts:}
\begin{parts}
  \item Define vapour pressure in terms of equilibrium.
        \answerlines[2]{The pressure exerted by vapour above a liquid in a
        closed container when the rate of evaporation equals the rate of
        condensation.}
  \item Does vapour pressure depend on how much liquid is present? Explain.
        \answerlines[2]{No --- it is set by temperature and intermolecular
        forces alone. A drop and a litre give the same value at the same
        temperature.}
  \item State the condition for boiling.
        \answerlines[2]{The liquid's vapour pressure equals the external
        pressure; the normal boiling point is where it reaches
        \SI{1}{\atm}.}
  \item Explain why water boils below \SI{100}{\celsius} at high altitude.
        \answerlines[2]{Atmospheric pressure is lower there, so the vapour
        pressure needed to match it is reached at a lower temperature.}
\end{parts}

\problem[]{Heating-curve segments. State whether each uses $q = mc\Delta T$
or $q = n\Delta H$, and explain what the energy does:}
\begin{parts}
  \item Warming ice from \SI{-20}{\celsius} to \SI{0}{\celsius}:
        \answerblank[6cm]{$mc\Delta T$ --- raises kinetic energy}
  \item Melting ice at \SI{0}{\celsius}:
        \answerblank[6cm]{$n\Delta H_{\text{fus}}$ --- overcomes IMFs}
  \item Boiling water at \SI{100}{\celsius}:
        \answerblank[6cm]{$n\Delta H_{\text{vap}}$ --- separates molecules}
\end{parts}

\problem[]{Calculations:}
\begin{parts}
  \item Energy to melt \SI{36.04}{\gram} of ice at \SI{0}{\celsius}:
        \answerblank[3.4cm]{\SI{12.0}{\kilo\joule}}
  \item Energy to vaporize \SI{36.04}{\gram} of water at
        \SI{100}{\celsius}: \answerblank[3.4cm]{\SI{81.4}{\kilo\joule}}
  \item Energy to warm \SI{36.04}{\gram} of liquid water from
        \SI{0}{\celsius} to \SI{100}{\celsius}: \workspace[1.8cm]
        \keyonly{\begin{apcanswer}$q = (36.04)(4.184)(100.0) =
        \SI{15.1}{\kilo\joule}$\end{apcanswer}}
  \item Which of these three steps requires the most energy, and why?
        \answerlines[2]{Vaporization --- separating molecules completely
        costs far more than loosening them enough to flow, or than simply
        raising their kinetic energy.}
\end{parts}

\problem[]{Explain why $\Delta H_{\text{vap}}$ (\SI{40.7}{\kilo\joule\per\mole})
is nearly seven times $\Delta H_{\text{fus}}$
(\SI{6.02}{\kilo\joule\per\mole}) for water.
\answerlines[3]{Melting only loosens the hydrogen-bonded lattice enough for
molecules to slide past one another --- most intermolecular contacts remain.
Vaporizing must break essentially all remaining attractions and move each
molecule far from its neighbours.}}

\problem[]{\textbf{FRQ (10 points).} A \SI{90.1}{\gram} sample of ice at
\SI{-15.0}{\celsius} is heated until it becomes steam at
\SI{115.0}{\celsius}.}
\begin{parts}
  \item Determine the number of moles of water.
        \answerlines[1]{$90.1/18.02 = \SI{5.00}{\mole}$}
  \item Calculate the energy required for each of the five stages.
        \workspace[4.2cm]
        \keyonly{\begin{apcanswer}Warm ice: $(90.1)(2.09)(15.0) =
        \SI{2.82}{\kilo\joule}$. Melt: $(5.00)(6.02) =
        \SI{30.1}{\kilo\joule}$. Warm water: $(90.1)(4.184)(100.0) =
        \SI{37.7}{\kilo\joule}$. Vaporize: $(5.00)(40.7) =
        \SI{204}{\kilo\joule}$. Warm steam: $(90.1)(2.0)(15.0) =
        \SI{2.7}{\kilo\joule}$.\end{apcanswer}}
  \item Determine the total energy required.
        \answerblank[3.4cm]{\SI{277}{\kilo\joule}}
  \item What percentage of the total is spent on vaporization? Comment on
        the result. \answerlines[2]{$204/277 = 74\%$. Nearly three-quarters
        of the energy goes into a single step that produces no temperature
        change at all.}
  \item Sketch the heating curve, labelling both plateaus and stating why
        the temperature is constant along them. \workspace[4cm]
        \keyonly{\begin{apcanswer}Curve rises, plateaus at
        \SI{0}{\celsius}, rises, plateaus (longer) at
        \SI{100}{\celsius}, rises. Temperature is constant because the
        energy is overcoming intermolecular forces rather than increasing
        average kinetic energy.\end{apcanswer}}
\end{parts}

\begin{apcnote}[Rubric --- score yourself, 10 points]
\textbf{(a) 1 pt} 5.00 mol \textbullet\ \textbf{(b) 4 pts} 1 per correct
sloped-segment pair (ice and steam), 1 fusion, 1 vaporization, 1 liquid
warming --- using mass in $n\Delta H$ or moles in $mc\Delta T$ loses the
point \textbullet\ \textbf{(c) 1 pt} \SI{277}{\kilo\joule} \textbullet\
\textbf{(d) 2 pts} 1 percentage, 1 comment noting no temperature change
\textbullet\ \textbf{(e) 2 pts} 1 correct shape with the longer plateau at
\SI{100}{\celsius}, 1 explanation citing IMFs versus kinetic energy.
\end{apcnote}

\problem[]{Explain why steam at \SI{100}{\celsius} produces far more severe
burns than liquid water at \SI{100}{\celsius}, using a quantitative
comparison. \answerlines[3]{Condensing steam releases
\SI{40.7}{\kilo\joule\per\mole} into the skin \emph{before} the resulting
water begins to cool. Cooling that same water from \SI{100}{\celsius} to
body temperature releases only about \SI{4.7}{\kilo\joule\per\mole}, so the
condensation delivers roughly nine times more energy.}}

\begin{spiralstrip}{Chapter 10 \textbullet\ blocks 1--4}
  \item Type of solid: diamond? \answerblank[3.4cm]{network covalent}
  \item Steel is which alloy type? \answerblank[2.6cm]{interstitial}
  \item Why does ice float? \answerblank[4.6cm]{open H-bonded lattice}
  \item Ionic solids conduct when: \answerblank[4.1cm]{molten or dissolved}
  \item Strongest IMF in \ce{NH3}: \answerblank[3.8cm]{hydrogen bonding}
\end{spiralstrip}

\end{document}
